\SourceLeaf{459}
\noindent Indem man diese beiden Gleichungen mit einander dividirt und dann die Quadratwurzel auszieht, ergiebt sich
\[
\frac{\sin\tfrac12 a^3}{\cos\tfrac12 a}.\frac{\cos\tfrac12 b}{\sin\tfrac12 b^3}
=\frac{\sin(A+\omega)\sin(A-\tfrac12\omega)^2}{\sin(B+\omega)\sin(B-\tfrac12\omega)^2}.
\]
\noindent Man kann diese Gleichung auch in die Form setzen
\[
\frac ab=\frac{\sin A}{\sin B}.\mkern6mu{}^3\mkern-4mu\BareRoot{D},
\]
\Needspace{10\baselineskip}
\noindent wo zur Abkürzung $D$ anstatt
\[
\begin{aligned}
&\frac{a^3\cos\tfrac12 a}{8\sin\tfrac12 a^3}.\frac{8\sin\tfrac12 b^3}{b^3\cos\tfrac12 b}\, .\\[5pt]
&\quad\frac{\sin(A+\omega)\sin(A-\tfrac12\omega)^2}{\sin A^3}.\frac{\sin B^3}{\sin(B+\omega)\sin(B-\tfrac12\omega)^2}
\end{aligned}
\]
\noindent geschrieben ist.

Diese Formel ist strenge richtig: man sieht aber sofort, dass wenn $a,b,c$ sehr klein sind, und als Grössen erster Ordnung betrachtet werden, jeder der vier Factoren, aus denen $D$ zusammengesetzt ist, von der Einheit nur um Grössen vierter Ordnung abweicht.

Nach allgemeinern Principien ist dieser Gegenstand abgehandelt und auf die Dreiecke ausgedehnt, die auf irgendwelchen krummen Flächen zwischen kürzesten Linien gebildet werden, in meinen \textit{Disquisitiones generales circa superficies curvas.}

